\section{Our Implementation: Components}
\label{sec:components}

The implementation is the ERC20 Vault contract.  Consult the \href{\analysedReadmeUrl}{ERC20 Vault README}  for a
high-level description of the contract logic and components.


\subsection{System Components}
\noindent
\textbf{Trusted Bridge $\SigSys$'s MPC.}
The instantiation of the cross-chain system relies on $\SigSys$'s MPC system acting as
a trusted bridge that listen for events (actions such as $\Deposit,\Swap,$ etc) on  Midnight
and executed them on Ethereum.
Between a request   (e.g., $\Deposit$) and its
settlement  circuit, three steps run outside any proof : the committee signs the requested EVM transaction;
a relayer broadcasts it; the committee observes the result and posts a signed attestation of it back
to Midnight.


$\SigSys$'s MPC  Signing System  is trusted to satisfy:
    \begin{enumerate}
        \item  \textbf{MPC committee integrity:}  the committee signs only what the protocol authorises.
        \item \textbf{Attestation faithfulness:}  MPC only provide attestations for events that have been executed on $\Lpub$ (Ethereum). The MPC attests failed only if the transaction did not execute.
    \end{enumerate}.

\begin{center}
\begin{tabular}{@{}L{2.7cm}L{5.4cm}L{6.3cm}@{}}
\toprule
\textbf{Definition} & \textbf{Instantiated by} & \textbf{Notes} \\
\midrule
$\Lpub$ &  Ethereum & accounts are EVM addresses; types are ERC20
  contracts; everything public \\
  \midrule
$\Lshd$ & Midnight  &  enables private payments  and execution of contract with private inputs.\\
\midrule
$\SigSys$  & SigNetwork MPC &  It manages the vault account on Ethereum and provides the attestations for it.\\
% \midrule
% $\Setup$ & deploy, $\cktInit$, $\cktApprove$ & \S\ref{ssec:setupinst} \\
% $\Deposit$ & $\cktDeposit$ $\to$ MPC $\to$ $\cktCompleteDeposit$ & \S\ref{ssec:depinst} \\
% $\Withdraw$ & $\cktWithdraw$ $\to$ MPC $\to$ $\cktCompleteWithdraw$ or $\cktRefund$ &
%   \S\ref{ssec:wdinst} \\
% $\Swap$ & $\cktSwap$ $\to$ MPC $\to$ $\cktCompleteSwap$ or $\cktRefund$ & \S\ref{ssec:swapinst} \\
\bottomrule
\end{tabular}
\end{center}


% $\Setup$ additionally publishes $\SigSys$'s public key ${\sf PK}_\SigSys$ in $\pp$.
% \medskip


\subsubsection{Key derivation}
\label{ssec:derive}

$\SigSys$ holds a single root key pair $(\mathsf{sk}_{\mathsf{root}}, \mpcRootPk)$ on
secp256k1, with $\mpcRootPk = \mathsf{sk}_{\mathsf{root}} \cdot G$ for the group generator
$G$ of order $n$. Every account it signs with is an \emph{additively derived} child of that
root, scoped by the requesting contract and by a $32$-byte path.


\begin{definition}[Additive key derivation]
\label{def:derive}
Let $\varepsilon : \{0,1\}^{*} \to \mathbb{Z}_n$ be
\[
  \Epsilon(\vaultAddr, \keyPath) \;=\;
  \Kec\bigl(\, \text{prefix} \,\|\, \chainId \,\|\, \vaultAddr \,\|\, \keyPath \,\bigr)
  \bmod n ,
\]
The derived key pair is
\begin{flalign*}
  &\mathsf{sk}_{\mathsf{der}} \;=\; \mathsf{sk}_{\mathsf{root}} + \varepsilon(\vaultAddr, \keyPath) \bmod n ,& \\
  &\Derive(\mpcRootPk\bc \vaultAddr\bc \keyPath) \;=\; \mpcRootPk + \Epsilon(\vaultAddr, \keyPath) \cdot G .&
\end{flalign*}
Where an EVM account is wanted, $\derivedAddr$ is the standard address encoding of the
derived public key: the low $20$ bytes of $\Kec$ applied to its uncompressed encoding.
\end{definition}

\noindent Three properties of Definition~\ref{def:derive} are used throughout.

\begin{itemize}[leftmargin=1.4em,itemsep=3pt,topsep=3pt]
\item \textbf{Public computability.} $\Derive$ takes only public inputs and addition on the
      curve is public, so \emph{anyone} can compute $\derivedAddr$ from $\mpcRootPk$,
      $\vaultAddr$ and $\keyPath$. The private half requires $\mathsf{sk}_{\mathsf{root}}$
      and therefore exists only inside $\SigSys$. This is what makes a deposit's
      $\keyPath = \userComm(\usersk)$ sufficient to expose the depositor's EVM address to an
      observer of $\Lshd$ alone.
\item \textbf{Contract scoping.} $\vaultAddr$ enters the derivation string, so no contract
      can reach another contract's derived accounts.
\item \textbf{Path separation.} Within one contract, distinct paths give unrelated keys.
      The vault uses exactly three: $\userComm(\usersk)$ for a depositor's account,
      the constant $\keyPathLit$ for the vault's own pooled account, and the fixed string
      $\text{``midnight response key''}$ for $\mpcRespKey$, which signs attestations and
      never signs transactions.
\end{itemize}

\noindent The key version $\keyVer$ selects which root key is used and is carried in every
request record; it is suppressed above, and throughout the document, since a single version
is in use.


\medskip
\noindent
\textbf{Derived Addresses:}

The value $\keyPath$ can be:
\begin{itemize}
    \item User address: $\derivedAddr$ uses $\keyPath=\userComm(\usersk)$ (defined below)
    \item Vault address:  $\vaultEvmAddr$  uses $\keyPath=\mathtt{''vault''}$
    \item  MPC Response Address: $\mpcRespKey$  uses $\keyPath=\mathtt{''midnight response key''}$. The attestation key $\mpcRespKey$ is derived the
same way with a fixed protocol string and pinned at $\cktInit$.
\end{itemize}


\subsection{User's Keys}
\label{ssec:userskeys}
A user  $U_i$ of  the cross-chain system holds:
\begin{itemize}
    \item { \bf Application secret key:} $\usersk_i$: a private key for the cross-chain system: $\usersk_i\leftarrow\{0,1\}^{\lambda}$. This is randomly chosen upon the first Deposit request.
      \item \textbf{Public account }$\pubAddr_i$  on $\Lpub$; (used to fund the derived account).
      \item \textbf{Shielded Account } ($\callerZKey_i$, $\callerZKeySk_i$) on $\Lshd$. \\
      In Midnight a shielded account consists of two public keys:
\begin{enumerate}
    \item \textbf{Coin Key-Pair} (Spending Key):
    \begin{enumerate}
        \item $\coinsk$ = $\SHA\bigl(\text{``midnight:csk''} \concat \seed\bigr)$
        \item $\coinpk = \SHA\bigl(\text{``midnight:zswap-pk[v1]''} \concat \coinsk\bigr)$
    \end{enumerate}

    \item \textbf{Encryption Key-Pair} (Discovery key):
    \begin{enumerate}
          \item $\encsk$: a group element generated via Poseidon hash function.
        \item $\encpk = \gen^{\encsk}$
    \end{enumerate}
\end{enumerate}
\noindent So
\[
  \shdAcct \;=\; \bigl(\,(\coinsk, \coinpk),\;\; (\encsk, \encpk)\,\bigr)
\]

\begin{center}
\begin{tabular}{@{}L{1.7cm}L{1.7cm}L{11.0cm}@{}}
\toprule
\textbf{Key} & \textbf{Exposure} & \textbf{Use} \\
\midrule
$\coinpk$ & public & \textbf{ownership}: used as a rcpt when Minting \\
$\coinsk$ & secret & \textbf{spending}:  used as sender when Spending \\
$\encpk$  & public & \textbf{payability}: used to compute ciphertext in Spending a coin\\
$\encsk$  & secret & \textbf{scanning} : used to decrypt in order to Receive a coin \\
\bottomrule
\end{tabular}
\end{center}
\end{itemize}


The  circuits  of our contract only use the $\coinpk$ as a parameter to mint shielded coins.
The other keys are used by the Midnight's circuits.


\subsubsection{User Authentication Values}
\label{ssec:derived}
\label{ssec:hashcomputation}
We use the following hash functions:
\begin{itemize}
    \item $\Hash$: A keyed hash function, used to compute our authentication tags. We require that this function satisfies Pseudorandomness (Def.~\ref{def:prf}) and Multi-user Pseudorandomness (Def.~\ref{def:muprf}) in addition to collision-resistance. Recall also that collision-resistance implies one-wayness (pre-image resistance).

    \begin{align*}
\hypertarget{fn:userComm}{\userComm(\usersk)} &= \Hash\bigl(\text{``vault:user:''}\bc \usersk\bigr)
  &&\text{\small identity; the MPC path in Deposit. }\\
\hypertarget{fn:refundComm}{\refundComm(\usersk, \rid)} &= \Hash\bigl(\text{``vault:refund:''}\bc \usersk\bc \rid\bigr)
  &&\text{\small the authentication tag.}
\end{align*}
We instantiate $\Hash$ with the \texttt{transientHash} implemented by Midnight~\footnote{\href{https://github.com/midnightntwrk/midnight-ledger/blob/cc909e042a17695f7ba62fc195a8d770d150b63e/transient-crypto/src/hash.rs}{Lines L75-L81}.}(based on Poseidon~\cite{poseidon}).
Poseidon follows a sponge-like approach and is a good candidate for Multi-user PRF security~\cite{KeyedDuplex17}.

    \item $\Kec$: A collision resistance hash function used to compute digests.
    \begin{align*}
\hypertarget{fn:domSep}{\domSep(\tokType)} &= \Hash\bigl(\text{``erc20:vault:''}\bc \tokType\bigr)
  &&\text{\small one shielded type per ERC20}\\
 \hypertarget{fn:reqIdOf}\rid &= \Kec(\req)
   &&\text{\small over the whole request record (\textbf{except skemas})}\\
%\hypertarget{fn:attDigest}{\attDigest(\rid, \evmOut)} &= \Kec(\rid \concat \evmOut)
%\newversion{}
\end{align*}
\end{itemize}


The \textbf{outcome} the MPC attests, is computed using the following formula:
 \begin{align*}
 \hypertarget{fn:attDigest}{\attDigest(\rid, \evmOut)}  = & \Kec(\rid \concat \evmOut.\kind \| \evmOut.\outlt \| \evmOut.\blockheight)
\end{align*}
  % &&\text{\small what the committee attests}
We instantiate $\Hash$ with the \texttt{transientHash} implemented by Midnight.

\subsection{The vault's public state}
\label{ssec:state}
\begingroup
\setlength{\LTpre}{6pt}\setlength{\LTpost}{6pt}
\begin{longtable}{@{}L{0.01cm}L{2.9cm}L{4.7cm}L{7.5cm}@{}}
\toprule
& \textbf{Here} & \textbf{Source name} & \textbf{Role} \\
\midrule
\endfirsthead
\multicolumn{4}{@{}l}{\small\itshape\ldots continued from previous page}\\
\toprule
& \textbf{Here} & \textbf{Source name} & \textbf{Information Held} \\
\midrule
\endhead
\midrule
\multicolumn{4}{r@{}}{\small\itshape continued on next page \ldots}\\
\endfoot
\bottomrule
\endlastfoot
 & $\reqMap$ & \srcReqMap & withdraw and approve requests, keyed by $\rid$; the MPC reads
  them here.  \\
& $\depositMap$ & \srcDepositMap & deposit requests keyed by $\rid$ \\
& $\swapMap$ & \srcSwapMap & swap requests keyed by $\rid$ \\
& $\supplyMap$ & \srcSupplyMap & supply requests keyed by $\rid$\\
& $\redeemMap$ & \srcRedeemMap & redeem requests  keyed by $\rid$\\
\addlinespace[3pt]
 & $\depositView$ & \srcDepositView & pending-\emph{deposit} record: the depositor
  commitment, the ERC20 and the amount $\cktCompleteDeposit$ mints \\
 & $\wdView$ & \srcWdView & pending-\emph{withdrawal} record: the refund commitment, the
  ERC20 and the amount.  \\
 & $\swapView$ & \srcSwapView & pending-\emph{swap} record: the refund commitment,
  $\tokType$, $\tokTypeB$, $\amtOut$ and $\amtMax$.  \\
& $\supplyView$ & \srcSupplyView & pending-\emph{supply} record: the refund commitment
  and the amount.\\
& $\redeemView$ & \srcRedeemView & pending-\emph{redeem} record: the refund commitment
  and the share count.  \\
\addlinespace[3pt]
& $\mpcRespKey$ & \code{mpcResponseKey} & signature verification key for attestation \\
& $\sigNonce$ & \code{signetRequestNonce} & a counter making otherwise identical requests distinct;
  its value publicly orders every request the vault has made \\
 & $\initFlag$ & \srcInitFlag & one-shot setup flag \\
& $\vaultEvmAddr$ & \code{vaultEvmAddress} & the contract-enforced recipient of every deposit \\
& $\routerAddr$ & \code{uniswapRouter} & the only spender $\cktApprove$ approves, and the only
  destination a swap request sends to \\
& $\underAddr$ & \srcUnderlying & the lent ERC20; gives $\cktStartSupply$ its burn color and
  $\cktCompleteRedeem$ its mint color \\
& $\wrapAddr$ & \srcWrapper & the interest-bearing wrapper; gives $\cktStartRedeem$ its burn color
  and $\cktCompleteSupply$ its mint color. The only spender $\cktApproveStata$ approves, and the
  only destination a supply or redeem request sends to \\
\end{longtable}
\endgroup

\noindent Also pinned at setup and public: the vault's own address $\vaultAddr$, the EVM chain
identifier, the Signet singleton, and the deployer's identity commitment


\vspace{5pt}
\newversion{The latest release replaces $\sigNonce$ and the five settle-view maps with a
request queue shared by every action: $\inputRequestBuffer$ and $\outputRequestBuffer$ hold
requests before and after nonce assignment, $\inputAttestationBuffer$ and
$\outputAttestationBuffer$ hold verified attestations before and after they are folded into
the block-height watermark $\lastseen$, $\vaultNonce$ is the vault EVM account's next nonce,
and \code{evictionMap} maps a sent $\rid$ back to its queue entry. Each action keeps a map of
its queued arguments (\code{depositArgsMap}, \code{withdrawArgsMap}, \ldots) and a map of its
sent requests (\code{bidirectionalDepositMap}, \ldots). A queue entry carries exactly one
user-dependent value, the ownership commitment
$\ownComm(\inIndex, \usersk) = \Hash(\usersk \concat \inIndex \concat \text{``vault:request:''})$
over a caller-chosen random index $\inIndex$, which stands in for $\refundComm(\usersk, \rid)$.
Everything else in these tables is bookkeeping (nonces, block heights, argument hashes and
attestation digests) that does not depend on the caller, so the tables do not alter the
privacy proof. They preserve the property that each request is unique and processed once:
completion consumes the entry, and a second identical request is refused while the first is
open.}
%-------------------------------------------------------------------------------------
%-------------------------------------------------------------------------------------
\subsection{The shielded Layer: Midnight Computation}
\label{ssec:shieldedcomputation}

Midnight provides shielded coin payments and private circuit evaluation.
We highlight with  \mnhl{shaded blue} the computation that are strictly Midnight's layer  and not the vault's own logic.

\vspace{5pt}
\noindent
\textbf{Midnight shielded-coin layer.}
In Midnight,  a raw secret coin is the tuple:
$c=(\mathsf{nonce}, \mathsf{color}, \mathsf{value})$.
Here are the functions applied to the secret coin:
\begin{align*}
\hypertarget{fn:coinCommit}{\coinCommit(c, \mathsf{rcpt})} &= \PersistantHash\bigl(\text{cc-tag} \concat \mathsf{nonce} \concat
  \mathsf{color} \concat \mathsf{value} \concat \mathsf{rcpt}\bigr)\\
\hypertarget{fn:coinNull}{\coinNull(c, \mathsf{sndr})} &= \PersistantHash\bigl(\text{cn-tag} \concat \mathsf{nonce} \concat
  \mathsf{color} \concat \mathsf{value} \concat \mathsf{sndr}\bigr)\\
\hypertarget{fn:tokenTypeOf}{\tokenTypeOf(\domSep(\tokType), \vaultAddr)} &= \PersistantHash\bigl(\text{tok-tag} \concat
  \domSep(\tokType) \concat \vaultAddr\bigr)
\end{align*}


% {\em User's coins vs contract coins.} For a \emph{user}-held coin the
% recipient  holds the coin \emph{public} key and the sender slot the coin \emph{secret} key, so
% an observer holding $\coinCommit$ cannot compute $\coinNull$. For a \emph{contract}-held coin both
% slots hold the same public contract address.


\medskip
\noindent\textbf{ Midnight Private-Circuit Computation.}
Midnight's contract computation is performed through circuits.
Every argument of a circuit is a \emph{private} input
of the proof.
A call to a  circuit $C$ on private arguments $\mathsf{in}$ and witness $\wit$ produces public outputs and state changes
for the vault contract $\Dvault$ and the Midnight chain $\Dshd$;
it also produces a statement $\stmt$ and a zk proof $\zkproof$.
\[
  \bigl(\; \Dvault,\;\; \Dshd,\;\; \stmt,\;\; \zkproof \;\bigr).
\]

\begin{center}
\begin{tabular}{@{}L{1.6cm}L{5cm}L{5cm}@{}}
\toprule
\textbf{Output} & \textbf{It contains} & \textbf{Chain Updates} \\
\midrule
$\Dvault$ & the list of assignments the call makes to the vault's public entries of
  \S\ref{ssec:state} & \textbf{the vault's} state on $\Lshd$ \\
\addlinespace[2pt]
$\Dshd$ & the list of coin commitments, nullifiers and mint records the call declares &
  \bluedot\textbf{Midnight's} state, not the vault's.\\
\addlinespace[2pt]
$\stmt$ & the public statement: the tuple an observer of $\Lshd$ reads off the call &
  --- \\
\addlinespace[2pt]
$\zkproof$ & a proof that $\Rel_C(\stmt, \wit)$ holds & --- \\
\bottomrule
\end{tabular}
\end{center}

 $\Dvault$, $\Dshd$ and $\stmt$  respect the following inclusion:

\[
  \Dvault \;\subseteq\; \stmt, \qquad \Dshd \;\subseteq\; \stmt,
  \qquad\text{and in general}\qquad \stmt \;\supsetneq\; \Dvault \cup \Dshd .
\]
\begin{itemize}[nosep,leftmargin=1.4em,topsep=3pt]
\item $\stmt$ also carries values that  do not update any state: the
      identifier $\rid$,  the commitment $\comIO$, the
      attestation the caller presented, and $\zkproof$ itself.
\item $\Rel_C$ is a \emph{predicate over $\stmt$ and $\wit$}, not a state update. This statement is verified inside the proof. Everything in $\stmt$ is visible.
\end{itemize}

\medskip
\noindent
\textbf{Midnight   Commitments.} We write $\comIO$ to denote the commitments  to the inputs and outputs of the circuit being called which are produced by the circuit computation.
The statement $\stmt$ contains $\comIO$.


\medskip
\noindent\textbf{How to read a circuit box.} \emph{Arguments} are the  inputs to the circuit (private);
\emph{witness} are the secret used to prove relations,   \emph{reads} are the public entries consulted; \emph{steps}
are the assertions and computations in source order;  A
\bluedot marked line or a blue block  highlights  Midnight's functions.


% \subsection{The attested outcome value}
% \label{ssec:outcomeval}

% Every settlement call consumes an attested outcome $\evmOut$, presented by the caller and
% authenticated only through $\attDigest(\rid, \evmOut)$. Exactly two kinds of value occur.

% \begin{definition}[$\succeeded$ and $\failed$]
% \label{def:evmout}
% Fix a request $\rid$ recorded by a request call. The attested outcome $\evmOut$ posted by $\SigSys$
% for $\rid$ is one of:
% \begin{itemize}[nosep,leftmargin=1.4em,topsep=2pt]
% \item a \emph{success value}, written $\evmOut \in \succeeded$: the encoding, under the response
%       schema $\schemaOut$ named in $\req$, of the return value of the EVM call that $\req$
%       requested. Its width and contents depend on the schema --- a one-byte boolean for a transfer,
%       an eight-byte integer for a swap, a supply or a redeem.
% \item the \emph{failure value}, written $\failed$: a single fixed protocol constant, independent of
%       the schema and of the request, posted when the requested transaction never executed. It is
%       the value $\failed$ that gates $\cktRefund$.
% \end{itemize}
% $\succeeded$ and $\failed$ are disjoint, and this disjointness is what routes an outcome to one
% settlement circuit rather than another.
% \end{definition}

%-------------------------------------------------------------------------------------
% ---------------------------------------------------------------------


\section{Our Implementation: Circuits}
\label{sec:instantiation}

In this section we show our implementation of the interface of Definition~\ref{def:system}.

\noindent
\textbf{Two calls per flow.}
Each procedure of Definition~\ref{def:system} is performed by the circuits below.
Each action: $\Deposit$, $\Withdraw$,  $\Swap$, $\Supply$, $\Redeem$, is performed by \emph{two} calls to the vault contract, with the $\SigSys$ bridge in
between. The first call records what the MPC should sign and takes
the values relevant to the call (transaction parameters). The second presents $\SigSys$'s attestation of what
happened on $\Lpub$ (Ethereum), and complete the request on $\Lshd$ (Midnight). Request calls are prefixed with \textsf{start}; settlement calls are prefixed with \textsf{complete}
The following table maps each call to their circuits.

\begin{center}
\begin{tabular}{@{}L{2.4cm}L{11.8cm}@{}}
\toprule
\textbf{Procedure} & \textbf{Circuits that implement it} \\
\midrule
$\Setup$    & $\cktInit$ (\cktref{ckt:initialize}); $\cktApprove$ (\cktref{ckt:approverouter}); $\cktApproveStata$ (\cktref{ckt:approvestata}) \\
$\Deposit$  & $\cktStartDeposit$ (\cktref{ckt:deposit}), then $\cktCompleteDeposit$ (\cktref{ckt:completedeposit}) \\
$\Withdraw$ & $\cktStartWithdraw$ (\cktref{ckt:startwithdraw}), then $\cktCompleteWithdraw$ (\cktref{ckt:completewithdraw}) or $\cktRefundWithdraw$ (\cktref{ckt:refundwithdraw}) \\
$\Swap$     & $\cktStartSwap$ (\cktref{ckt:startswap}), then $\cktCompleteSwap$ (\cktref{ckt:completeswap}) or $\cktRefundSwap$ (\cktref{ckt:refundswap}) \\
$\Supply$   & $\cktStartSupply$ (\cktref{ckt:startsupply}), then $\cktCompleteSupply$ (\cktref{ckt:completesupply}) or $\cktRefundSupply$ (\cktref{ckt:refundsupply}) \\
$\Redeem$   & $\cktStartRedeem$ (\cktref{ckt:startredeem}), then $\cktCompleteRedeem$ (\cktref{ckt:completeredeem}) or $\cktRefundRedeem$ (\cktref{ckt:refundredeem}) \\
\bottomrule
\end{tabular}
\end{center}


\newversion{In the latest release the global counter $\sigNonce$ is dropped. A request id
is made distinct by the EVM transaction nonce instead, and freshness comes from the
block-height watermark $\lastseen$. Transaction nonces are not assigned at request time by
the caller of {\code{startX}} (for any action X except Deposit, whose caller names the nonce
of their own derived account), but by a dedicated circuit that any party can call.

Specifically, each circuit {\code{startX}} only validates and records the transaction
parameters, without the EVM nonce and without the request id, under a caller-chosen random
index $\inIndex$. This \emph{partial} request is added to the buffer $\inputRequestBuffer$ of
requests waiting for nonce assignment. Every so often the queue is processed by the circuit
$\cktFlush$, which assigns the vault account's next nonce $\vaultNonce$ (a counter only
$\cktFlush$ reads or writes) to each vault-signed request and stamps every request with the
current $\lastseen$, and by the per-action circuits $\cktSend$ (\code{sendDeposit},
\code{sendWithdraw}, \ldots), which build the complete request from the flushed entry, compute
$\rid = \reqIdOf(\req)$ and notify the MPC. On the way back, the circuits $\cktQueueAtt$
verify the MPC's attestation signature and queue the attestation, $\cktFlush$ folds it into
$\outputAttestationBuffer$ and raises $\lastseen$ to its block height, and {\code{completeX}}
consumes the request together with its attestation.
$\cktFlush$, $\cktSend$ and $\cktQueueAtt$ are permissionless, take no witness and can be
executed by everyone.
Note, these circuits only add nonce and time-stamping information that is not user-dependent.
As such they do not alter the privacy analysis described in Sec.~\ref{sec:proof}. They
implement efficiency optimisations, contention avoidance and replay-attack protection. The
refund circuits ({\code{refundX}}) are folded into {\code{completeX}}, which refunds on a
failed or unviable verdict, and every {\code{completeX}}, the success branch of
$\cktCompleteWithdraw$ included, is gated on the requester's ownership commitment.
}


\medskip
\noindent
\textbf{Notification to $\SigSys$.} Every request call ends with a
cross-contract call to the Signet singleton, which emits an event the MPC polls. We write
$\notif(\vaultAddr, \rid, M)$ for that event: it carries the vault's address, the request
identifier, and the identity of the ledger map $M$ in which the request record was stored, so
that the MPC knows where to read $\req$ from.


\subsection{Setup: \texorpdfstring{$\Setup^{\Lpub,\Lshd}(1^{\secpar})$}{Setup}}
\label{ssec:setupinst2}

\begin{ProcBox}{$\pp \gets \Setup^{\Lpub,\Lshd}(1^{\secpar})$}
\sloppy Run once by the deployer, who holds a secret $\usersk_D$ of their own.
\begin{Steps}
\item \textbf{Deploy the contract} on $\Lshd$ with two constructor arguments: the deployer's
      identity $\deplComm = \lnk{fn:userComm}{\userComm(\usersk_D)}$ and the address of the Signet
      singleton. The contract's own address $\vaultAddr$ becomes known here, which is why the
      remaining values cannot be pinned in the constructor.
\item \textbf{Derive off-chain}
      \[\begin{aligned}
        \vaultEvmAddr &\gets \Derive(\mpcRootPk\bc \vaultAddr\bc \keyPathLit),\\
        \mpcRespKey &\gets \Derive(\mpcRootPk\bc \vaultAddr\bc \text{``midnight response key''}).
      \end{aligned}\]
\item \textbf{Run $\cktInit$} --- \cktref{ckt:initialize}.
\item \textbf{Run $\cktApprove$} once per ERC20 that will ever be swapped ---
      \cktref{ckt:approverouter}.
\end{Steps}
Output the public parameters:\\
$\pp = (\vaultAddr, \vaultEvmAddr, \mpcRespKey, \chainId, \routerAddr, \text{singleton address})$.
\end{ProcBox}


\begin{CktBox}{ckt:initialize}{$\cktInit$}
\fld{Arguments.} $\vaultEvmAddr$, $\routerAddr$, \ADDEDINL{$\underAddr$, $\wrapAddr$},
$\chainId$, $\mpcRespKey$.
\fld{Witness.} $\usersk$.
\fld{Reads.} $\initFlag$, $\deplComm$.
\fld{Steps.}
\begin{enumerate}[nosep,leftmargin=1.7em,topsep=2pt]
\item $\assert \initFlag = 0$.
\item $\assert \lnk{fn:userComm}{\userComm(\usersk)} = \deplComm$.
\item $\assert \chainId > 0$, \; $\routerAddr \neq 0$,
      \ADDEDINL{\; $\underAddr \neq 0$, \; $\wrapAddr \neq 0$}.
\item Write $\initFlag \gets 1$ and the six pinned values.
\end{enumerate}
\begin{OutList}
\oitem{\Dvault} $\initFlag \gets 1$; $\vaultEvmAddr$, $\routerAddr$,
  \ADDEDINL{$\underAddr$, $\wrapAddr$}, $\chainId$, $\mpcRespKey$ written.
\oitem{\Dshd} $\emptyset$.
\oitem{\stmt} $(\cktInit,\; \Dvault,\; \comIO)$.
\oitem{\wit} $\usersk$.
\oitem{\Rel} $\userComm(\usersk) =\deplComm$.
\end{OutList}
\end{CktBox}


\begin{CktBox}{ckt:approverouter}{$\cktApprove$}
\fld{Arguments.} $\tokType$, $\evmNonce$, $\keyVer$.
\fld{Witness.} none.
\fld{Reads.} $\initFlag$, $\routerAddr$, $\vaultAddr$, $\sigNonce$, $\chainId$.
\fld{Steps.}
\begin{enumerate}[nosep,leftmargin=1.7em,topsep=2pt]
\item $\assert \initFlag \geq 1$ and $\tokType \neq 0$.
\item $\evmTxOf \gets$ ``approve $\routerAddr$ for an unlimited amount of $\tokType$''. \emph{The
      spender and the amount come from pinned state; the caller chooses only the token.}
\item $\req \gets (\vaultAddr\bc \sigNonce\bc \keyVer\bc \keyPath{=}\keyPathLit\bc \evmTxOf\bc
      \chainId)$; \; $\rid \gets \lnk{fn:reqIdOf}{\reqIdOf(\req)}$; \;
      $\assert \rid \notin \reqMap$.
\end{enumerate}
\begin{OutList}
\oitem{\Dvault} $\sigNonce \mathrel{+}= 1$; \; $\reqMap[\rid] \gets \req$ ---
  \emph{never removed by any circuit}.
\oitem{\Dshd} $\emptyset$.
\oitem{\stmt} $(\cktApprove,\; \req,\; \rid,\; \notif(\vaultAddr, \rid, \reqMap))$.
\oitem{\wit} $\bot$.
\oitem{\Rel} $\mathsf{true}$ - the witness is empty and every argument is in $\stmt$, so
  $\zkproof$ certifies a computation an observer can redo.
\end{OutList}
\end{CktBox}

% ---------------------------------------------------------------------
\subsection{Deposit}
\label{ssec:depinst2}

\begin{ProcBox}{$\Pi.\Deposit^{\Lpub, \Lshd}\bigl([\pubAddr]\bc \shdAcct\bc \tokType\bc \amt\bc
      \auxx \,;\, (\secret{\usersk}, \freshrn{\mintNonce})\bigr)$.}
\sloppy
\emph{Private input.} secret key $\usersk$; and a fresh $\mintNonce$, used at Step~4.\\
\emph{Public arguments.} $\tokType$ is an ERC20 contract,
and $\auxx = (\evmNonce, \gasEnv, \keyVer)$.
The receiving \zkswap{} account $\shdAcct$  is not public per-se, but we
do not aim to protect it.

\begin{Steps}
\item \textbf{Fund the derived account:} outside any circuit, on $\Lpub$. The caller pays $\amt$ of $\tokType$  to  the MPC controlled account $\derivedAddr$ on $\Lpub$. Where,
      $\derivedAddr = \Derive\bigl(\mpcRootPk\bc \vaultAddr\bc
      \lnk{fn:userComm}{\userComm(\usersk)}\bigr)$. See \S\ref{ssec:derived}.
\item \textbf{Request call: $\cktDeposit$} (\cktref{ckt:deposit}). Recomputes
      $\keyPath \gets \lnk{fn:userComm}\userComm(\usersk)$ from the witness and request
      $\text{transfer}(\vaultEvmAddr, \amt)$ under $\rid$. \emph{$\Dshd = \emptyset$: no coin
      moves.}
\item \textbf{MPC}. The committee signs on behalf of  $\derivedAddr$ a transfer to $\vaultEvmAddr$, and posts $\attSig$ over
      $\lnk{fn:attDigest}{\attDigest(\rid, \evmOut)}$.
\item \textbf{Settlement call: $\cktCompleteDeposit$} (\cktref{ckt:completedeposit}). Verifies
      $\attSig$, consumes $\reqMap[\rid]$, authenticated the request on $\userComm(\usersk) = \req.\keyPath$, and mints
      $\amt$ to $\shdAcct$ using randomness $\mintNonce$.
\end{Steps}

\medskip
Outcome: $\succeeded$ once Step~4 completes; $\failed$ if the EVM transfer did not execute, in which
case the value never left $\derivedAddr$;
Private output: $\usersk,\mintNonce$.
\end{ProcBox}


\begin{CktBox}{ckt:deposit}{$\srcStartDeposit$}
\fld{Arguments.} $\evmNonce$, $\gasEnv$, $\keyVer$, $(\tokType, \amt)$.
\fld{Witness.} $\usersk$.
\fld{Reads.} $\initFlag$, $\vaultEvmAddr$, $\vaultAddr$, $\sigNonce$, $\chainId$.
\fld{Steps.}
\begin{enumerate}[nosep,leftmargin=1.7em,topsep=2pt]
\item $\assert \initFlag \geq 1$, \; $\tokType \neq 0$, \; $\gasEnv > 0$, \;
      $0 < \amt \leq 2^{64}-1$.
\item $\keyPath \gets \lnk{fn:userComm}{\userComm(\usersk)}$.
\item $\evmTxOf \gets$ ``$\text{transfer}(\vaultEvmAddr, \amt)$ on $\tokType$'', with $\evmNonce$,
      $\gasEnv$, on $\chainId$.
\item $\req \gets (\vaultAddr\bc \sigNonce\bc \keyVer\bc \keyPath\bc \evmTxOf\bc \chainId)$; \;
      $\rid \gets \lnk{fn:reqIdOf}{\reqIdOf(\req)}$; \;
      \ADDEDINL{$\assert \rid \notin \depositMap$}.
\end{enumerate}
\begin{OutList}
\oitem{\Dvault} $\sigNonce \mathrel{+}= 1$; \; \ADDEDINL{$\depositMap[\rid] \gets \req$}; \;
  \ADDEDINL{$\depositView[\rid] \gets (\keyPath,\; \tokType,\; \amt)$}.
\oitem{\Dshd} $\emptyset$.
\oitem{\stmt} $(\cktStartDeposit,\; \req,\; \rid,\; \depositView[\rid],\;
  \notif(\vaultAddr, \rid, \depositMap),\; \comIO)$, where $\req$ exposes
  $\tokType, \amt, \vaultEvmAddr, \evmNonce, \gasEnv, \sigNonce$ and
  $\keyPath = \userComm(\usersk)$.
\oitem{\wit} $\usersk$.
\oitem{\Rel} $\req.\keyPath = \userComm(\usersk)$ \;$\wedge$\;
  $\depositView[\rid].\fComm = \userComm(\usersk)$.
\end{OutList}
\end{CktBox}


\begin{CktBox}{ckt:completedeposit}{$\cktCompleteDeposit$}
\fld{Arguments.} $\rid$, $\attSig$, $\evmOut$, $\freshrn{\mintNonce}$,  $\shdAcct$.
\fld{Witness.} $\usersk$.
\fld{Reads.} $\initFlag$, $\mpcRespKey$, \ADDEDINL{$\depositMap$, $\depositView[\rid]$}.
\fld{Steps.}
\begin{enumerate}[nosep,leftmargin=1.7em,topsep=2pt]
\item $\assert \initFlag \geq 1$ and $\evmOut$ reports success.
\item $\assert \EcdsaVf\bigl(\lnk{fn:attDigest}{\attDigest(\rid, \evmOut)}\bc \attSig\bc
      \mpcRespKey\bigr)$ --- \textbf{gate 1}, the authentication of the MPC.
\item \ADDEDINL{$\assert \rid \in \depositMap$; \; remove $\depositMap[\rid]$} ---
      \textbf{gate 3}, at most once, enforced by the ledger through $\Dvault$, not by $\Rel$.
\item \ADDEDINL{$\assert \lnk{fn:userComm}{\userComm(\usersk)} = \depositView[\rid].\fComm$;} \;
      remove $\depositView[\rid]$ --- \textbf{gate 2}, enforces depositor-only.
\item \ADDEDINL{$\amt \gets \depositView[\rid].\fAmt$, \;
      $\tokType \gets \depositView[\rid].\fErc$}.
\end{enumerate}
\MidNote{\textbf{Midnight operation.} Execute Midnight Circuit $\mintShieldedToken$ on input
$\domSep, \amt, \mintNonce, \callerZKey$. Let $\coinMint =
\bigl(\;\freshrn{\mintNonce}, \mathsf{color}, \amt \bigr)$}
where  $\mathsf{color}=\lnk{fn:tokenTypeOf}{\tokenTypeOf}\bigl(\domSep(\tokType),
\vaultAddr\bigr)$
\begin{OutList}
\oitem{\Dvault} \ADDEDINL{$\depositMap[\rid]$ removed; $\depositView[\rid]$ removed}.
\oitem{\Dshd} \bluedot mint record $\bigl(\lnk{fn:domSep}{\domSep(\tokType)},\, \amt\bigr)$;
  \bluedot output  $\lnk{fn:coinCommit}{\coinCommit(\coinMint, \shdAcct)}$.
\oitem{\stmt} $\bigl(\cktCompleteDeposit,\; \rid,\; \attSig,\;
  (\domSep(\tokType), \amt),\; \coinCommit(\coinMint, \callerZKey)\bigr)$.
\oitem{\wit} $(\usersk,\; \mintNonce,\; \shdAcct)$.
\oitem{\Rel} $\userComm(\usersk)$ $= \depositView[\rid].\fComm$ \;$\wedge$\;
  $\EcdsaVf(\attDigest(\rid,\evmOut)$, $\attSig$, $\mpcRespKey)$ \;$\wedge$\;
  the published commitment opens to $(\mintNonce, \cdot, \amt, \callerZKey)$.
\end{OutList}
\end{CktBox}


\begin{PubBox}{Transcript from  $\Deposit$}
\textbf{On $\Lpub$:} [$\pubAddr \to \derivedAddr$],  and
$\derivedAddr \to \vaultEvmAddr$ (signed by the committee), both in the clear.

\smallskip
\textbf{On $\Lshd$:} $\stmt_{\cktDeposit}$ and
$\stmt_{\cktCompleteDeposit}$.

\smallskip
\textbf{An observer therefore holds} $\userComm(\usersk)$, hence $\derivedAddr$, hence the edge from
the caller's own EVM account to their identity in this system; $\tokType$ and $\amt$; $\rid$, so \textbf{the settlement call is  linked
to the request call}. Hidden: $\usersk$, $\mintNonce$.
\end{PubBox}
% ---------------------------------------------------------------------
\subsection{Withdraw}
\label{ssec:wdinst2}

\begin{ProcBox}{$\Withdraw^{\Lpub, \Lshd}\bigl(\shdAcct\bc \destAddr*\bc \tokType\bc \amt\bc
      \auxx \,;\, \privIn\bigr)$}\label{box:withdraw}
\sloppy
\emph{Private input.} $\privIn$ holds $\usersk$; the secret key $\coinsk$ associated to the $\zkswap$ account $\srcAcct$; a
fresh $\burnNonce$, used at Step~1; and, only if the flow ends in $\failed$, a fresh $\mintNonce$
used at Step~4.
\emph{Public arguments.} $\auxx = (\evmNonce, \keyVer)$ only: the gas envelope is fixed in the
contract, because the \emph{vault} pays for this transaction.

 \begin{Steps}
 \item \textbf{Request call: $\cktWithdraw$} (\cktref{ckt:startwithdraw}).
 The caller hands the contract  a coin descriptor   $\coinIn = \bigl(\burnNonce,\;\mathsf{color}, \amt\bigr)$; the circuit burns it, records  $\text{transfer}(\destAddr^*, \amt)$ under $\rid$ with $\keyPath = \keyPathLit$ o, and add the authentication tag to the map:  $\wdView[\rid] \gets \lnk{fn:refundComm}{\refundComm(\usersk, \rid)}$.
   \smallskip

 \item \textbf{MPC}. The MPC signs the evm request with the key derived
       at $\keyPath = \keyPathLit$, so from $\vaultEvmAddr$, the pool's own account. Once the request is executed,  the MPC signs an attestation $\attSig$ of the outcome $\evmOut$.
\item \textbf{Settlement call on success: $\cktCompleteWithdraw$}, success branch  (\cktref{ckt:completewithdraw}). \textbf{$\Dshd = \emptyset$ and $\wit = \bot$ on this branch}.
\item \textbf{Settlement call on failure: $\cktCompleteWithdraw$'s failure branch, or $\cktRefundWithdraw$}
      -  \cktref{ckt:completewithdraw} and \cktref{ckt:refundwithdraw}, both which gates on
      $\lnk{fn:refundComm}{\refundComm(\usersk, \rid)}$ against the stored marker  and re-mints $\amt$ under a fresh
      $\mintNonce$. The two circuits differ only in which attested output routes to them: an
      executed transfer that returned false goes to $\cktCompleteWithdraw$; a transaction that never
       executed goes to $\cktRefundWithdraw$.
 \end{Steps}

\medskip
Outcome: $\succeeded$ via Step~3, $\failed$ via Step~4. Private output: $\usersk$.
\end{ProcBox}


\begin{CktBox}{ckt:startwithdraw}{$\cktStartWithdraw$}
\fld{Arguments.} $\evmNonce$, $\keyVer$, $(\tokType, \amt, \destAddr)$, and
$\coinIn = (\burnNonce, \mathsf{color}, \mathsf{value})$.
\fld{Witness.} $\usersk$ (and $\coinsk$ at the wallet layer).
\fld{Reads.} $\initFlag$, $\vaultAddr$, $\sigNonce$, $\chainId$.
\fld{Steps.}
\begin{enumerate}[nosep,leftmargin=1.7em,topsep=2pt]
\item $\assert \initFlag \geq 1$, \; $\tokType \neq 0$, \; $0 < \amt \leq 2^{64}-1$.
\item $\assert \coinIn.\mathsf{color} =
      \lnk{fn:tokenTypeOf}{\tokenTypeOf(\domSep(\tokType), \vaultAddr)}$ \; and \;
      $\coinIn.\mathsf{value} = \amt$.
\item $\evmTxOf \gets$ ``$\text{transfer}(\destAddr, \amt)$ on $\tokType$'', with $\evmNonce$ and a
      contract-fixed gas envelope, on $\chainId$.
\item $\keyPath \gets \keyPathLit$.
\item $\req \gets (\vaultAddr\bc \sigNonce\bc \keyVer\bc \keyPath\bc \evmTxOf\bc \chainId)$; \;
      $\rid \gets \lnk{fn:reqIdOf}{\reqIdOf(\req)}$; \; $\assert \rid \notin \reqMap$.
\end{enumerate}
\MidNote{\textbf{Midnight operations: burn.}
$\mnhl{\code{receiveShielded}}(\coinIn)$ then
$\mnhl{\code{sendImmediateShielded}}(\coinIn, \burnAddr, \amt)$.}
\begin{OutList}
\oitem{\Dvault} $\sigNonce \mathrel{+}= 1$; \; $\reqMap[\rid] \gets \req$; \;
  \ADDEDINL{$\wdView[\rid] \gets \bigl(\lnk{fn:refundComm}{\refundComm(\usersk, \rid)},\;
  \tokType,\; \amt\bigr)$}.
\oitem{\Dshd} \bluedot $\coinCommit(\coinIn, \vaultAddr)$, \;
  \bluedot $\coinNull(\coinIn, \vaultAddr)$, \;
  \bluedot $\coinCommit(\coinBurnOut, \burnAddr)$, \;
  \bluedot $\{\coinNull(\coinFund_i, \coinsk)\}_i$ and $\coinCommit(\coinChange, \srcAcct)$.
\oitem{\stmt} $\bigl(\cktStartWithdraw,\; \req,\; \rid,\; \wdView[\rid],\; \Dshd,\;
  \notif(\vaultAddr, \rid, \reqMap),\; \comIO\bigr)$.
  \ADDEDINL{The view publishes $\tokType$ and $\amt$.}
\oitem{\wit} $(\usersk,\; \burnNonce,\; \coinsk)$.
\oitem{\Rel} $\coinIn.\mathsf{color} = \tokenTypeOf(\domSep(\tokType), \vaultAddr)$ \;$\wedge$\;
  $\coinIn.\mathsf{value} = \amt$ \;$\wedge$\;
  $\wdView[\rid].\fComm = \refundComm(\usersk, \rid)$.
\end{OutList}
\end{CktBox}


\begin{CktBox}{ckt:failgate}{$\cktFailGate$ - shared across refund circuits (\srcFailGate)}
Not exported. Called first by $\cktRefundWithdraw$, $\cktRefundSwap$, $\cktRefundSupply$ and
$\cktRefundRedeem$.
\fld{Arguments.} $\rid$, $\attSig$, $\evmOut$ of width $5$.
\fld{Reads.} $\initFlag$, $\mpcRespKey$.
\fld{Steps.}
\begin{enumerate}[nosep,leftmargin=1.7em,topsep=2pt]
\item $\assert \initFlag \geq 1$.
\item $\assert \EcdsaVf\bigl(\lnk{fn:attDigest}{\attDigest(\rid, \evmOut)}\bc \attSig\bc
      \mpcRespKey\bigr)$.
\item $\assert \evmOut = \failed$.
\end{enumerate}
\begin{OutList}
\oitem{\Dvault} none of its own.
\oitem{\Dshd} $\emptyset$.
\oitem{\stmt} $(\rid,\; \attSig,\; \evmOut)$, merged into the caller circuit's statement.
\oitem{\wit} $\bot$.
\oitem{\Rel} $\EcdsaVf(\attDigest(\rid,\evmOut), \attSig, \mpcRespKey)$ \;$\wedge$\;
  $\evmOut = \failed$.
\end{OutList}
\end{CktBox}


\begin{CktBox}{ckt:completewithdraw}{$\cktCompleteWithdraw$}
\fld{Arguments.} $\rid$, $\attSig$, $\evmOut$ of width $1$, $\mintNonce$.
\fld{Witness.} $\usersk$ - \emph{failure branch only}.
\fld{Reads.} $\initFlag$, $\mpcRespKey$, \ADDEDINL{$\wdView[\rid]$}, $\reqMap$.
\fld{Steps.}
\begin{enumerate}[nosep,leftmargin=1.7em,topsep=2pt]
\item $\assert \initFlag \geq 1$ and
      $\assert \EcdsaVf\bigl(\lnk{fn:attDigest}{\attDigest(\rid, \evmOut)}\bc \attSig\bc
      \mpcRespKey\bigr)$.
\item \ADDEDINL{$\assert \rid \in \wdView$}.
\item Remove $\reqMap[\rid]$. \ADDEDINL{The record is not read.}
\item $\okbit \gets$ the success bit of $\evmOut$.
\item \ADDEDINL{If $\okbit = 0$:
      $\assert \lnk{fn:refundComm}{\refundComm(\usersk, \rid)}$
      $= \wdView[\rid].\fComm$, then mint
      $\coinMint = \bigl(\mintNonce,\, \tokenTypeOf(\domSep(\wdView[\rid].\fErc), \vaultAddr),\,
      \wdView[\rid].\fAmt\bigr)$
      to the caller's own shielded account.}
\item Remove $\wdView[\rid]$.
\end{enumerate}
\begin{OutList}
\oitem{\Dvault} $\reqMap[\rid]$ removed; \ADDEDINL{$\wdView[\rid]$ removed}.
\oitem{\Dshd} $\okbit = 1$: $\emptyset$. \quad $\okbit = 0$: \bluedot mint record
  $\bigl(\domSep(\wdView[\rid].\fErc),\, \wdView[\rid].\fAmt\bigr)$; \;
  \bluedot $\coinCommit(\coinMint, \text{caller})$.
\oitem{\stmt} $\bigl(\cktCompleteWithdraw,\; \rid,\; \attSig,\; \okbit,\; \Dshd\bigr)$.
\oitem{\wit} $\okbit = 1$: $\bot$. \quad $\okbit = 0$: $(\usersk, \mintNonce)$.
\oitem{\Rel} $\okbit = 1$: $\EcdsaVf(\attDigest(\rid,\evmOut), \attSig, \mpcRespKey)$. \quad
  $\okbit = 0$: the same, $\wedge$ $\refundComm(\usersk, \rid)$ $= \wdView[\rid].\fComm$ $\wedge$
  the published commitment opens to $(\mintNonce, \cdot, \wdView[\rid].\fAmt, \text{caller})$.
\end{OutList}
\fld{$\wit = \bot$ on the success branch.} The proof certifies a statement with no secret in it, so
any third party holding the attestation can produce the same call. $\mintNonce$ is an argument on
both branches; on the success branch it is simply unused.
\end{CktBox}


\begin{CktBox}{ckt:refundwithdraw}{ $\cktRefundWithdraw$}
\fld{Arguments.} $\rid$, $\attSig$, $\evmOut$ of width $5$, $\mintNonce$.
\fld{Witness.} $\usersk$.
\fld{Reads.} $\initFlag$, $\mpcRespKey$, $\wdView[\rid]$, $\reqMap$.
\fld{Steps.}
\begin{enumerate}[nosep,leftmargin=1.7em,topsep=2pt]
\item Run $\cktFailGate(\rid, \attSig, \evmOut)$.
\item $\assert \rid \in \wdView$.
\item $\assert \lnk{fn:refundComm}{\refundComm(\usersk, \rid)} = \wdView[\rid].\fComm$ (authenticating the withdrawer).
\item Remove $\reqMap[\rid]$ and $\wdView[\rid]$.
\end{enumerate}
\MidNote{\textbf{Midnight operation.} Mint, to the caller's own shielded account,
\[\coinMint = \bigl(\mintNonce,\;
\lnk{fn:tokenTypeOf}{\tokenTypeOf(\domSep(\wdView[\rid].\fErc), \vaultAddr)},\;
\wdView[\rid].\fAmt\bigr).\]}
\begin{OutList}
\oitem{\Dvault} $\reqMap[\rid]$ removed; $\wdView[\rid]$ removed.
\oitem{\Dshd} \bluedot mint record
  $\bigl(\domSep(\wdView[\rid].\fErc),\, \wdView[\rid].\fAmt\bigr)$; \;
  \bluedot $\lnk{fn:coinCommit}{\coinCommit(\coinMint, \text{caller})}$.
\oitem{\stmt} $\bigl(\cktRefundWithdraw,\; \rid,\; \attSig,\; \Dvault,\; \Dshd\bigr)$.
\oitem{\wit} $(\usersk,\; \mintNonce)$.
\oitem{\Rel} $\refundComm(\usersk, \rid)$ $= \wdView[\rid].\fComm$ \;$\wedge$\; the relation of
  $\cktFailGate$ \;$\wedge$\; the published commitment opens to
  $(\mintNonce, \cdot, \wdView[\rid].\fAmt, \text{caller})$.
\end{OutList}
\end{CktBox}


\begin{PubBox}{What $\Withdraw$ publishes}
\textbf{On $\Lpub$:} on success, one transfer $\vaultEvmAddr \to \destAddr$ of $\amt$ of $\tokType$.
\emph{Nothing tied to $\srcAcct$ appears.}

\smallskip
\textbf{On $\Lshd$:} $\stmt_{\cktWithdraw}$, the settlement
statement. Of these, the components that involve the caller at all are: the marker
$\refundComm(\usersk, \rid)$;  the coins sent to the burn address $\burnAddr$ in the $\cktStartWithdraw$ (the identity of the payer for this coins is protected by Mindnight's library under the assumption that $\burnNonce$ is private and is random), and -- in case of failure -- the new coins minted to the caller of $\cktRefundWithdraw$ (the identity of the payer for this coins is protected by Mindnight's library under the assumption that $\mintNonce$ is private and is random).


\smallskip
\textbf{Hidden:} which shielded account funded the call/ asked for a refund.
\end{PubBox}


%---------------------------------------------------------------------
%---------------------------------------------------------------------------


\subsection{Swap}
\label{ssec:swapinst2}

\begin{ProcBox}{$\Pi.\Swap^{\Lpub, \Lshd}\bigl(\shdAcct\bc \tokType\bc \amt\bc
    \tokTypeB\bc \amtB\bc \auxx \,;\, \privIn\bigr)$}
\sloppy
\emph{Private input.} As $\Withdraw$: $\usersk$, $\coinsk$ of $\srcAcct$, a fresh $\burnNonce$ at
Step~1, and a fresh $\mintNonce$ at Step~3 and a fresh $\changeNonce$ in Step 4.).
\emph{Public arguments.} $\amt$ is a \emph{maximum} input $\amtMax$ and $\amtB$ an exact output
$\amtOut$; (in an honest use $\srcAcct$ is the recipient of the swapped coin);
$\auxx = (\evmNonce, \keyVer, \feeTier)$.
\begin{Steps}
\item \textbf{Request call: $\cktSwap$} - \cktref{ckt:startswap}. Same shape as $\cktWithdraw$: a   coin descriptor for $\amtMax$ of $\tokType$ is surrendered and burned, an exact-output exchange is
      recorded under $\rid$ in the map $\swapMap$, and
      $\swapView[\rid] \gets \lnk{fn:refundComm}{\refundComm(\usersk, \rid)}$ is recorded.
\item \textbf{MPC}. As in withdraw,   the MPC signs the swap request on behalf of  $\vaultEvmAddr$. Once the request is executed,  the MPC signs an attestation $\attSig$ of the outcome $\evmOut$.
\item \textbf{Settlement call on success: $\cktCompleteSwap$} - \cktref{ckt:completeswap}.
If the transaction is successful, the user can claim the new-colored coin and the change, after proving knowledge of the pre-image of $\refundComm$.

\item \textbf{Settlement call on failure: $\cktRefundSwap$}. If the transaction reverts on $\Lpub$, user calls  Circuit \cktref{ckt:refundswap} to re-mint a coin with the full $\amtMax$ of the original color, after proving  knowledge of the pre-image of $\refundComm$.
\end{Steps}

\medskip
Outcome: $\succeeded$ via Step~3, $\failed$ via Step~4.
Private output: the new coins minted.
\end{ProcBox}


\begin{CktBox}{ckt:startswap}{$\cktStartSwap$}
\fld{Arguments.} $\evmNonce$, $\keyVer$, $(\tokType, \tokTypeB, \feeTier, \amtOut, \amtMax)$, and
$\coinIn = (\burnNonce, \mathsf{color}, \mathsf{value})$.
\fld{Witness.} $\usersk$ (and $\coinsk$ at the wallet layer).
\fld{Reads.} $\initFlag$, $\routerAddr$, $\vaultEvmAddr$, $\vaultAddr$, $\sigNonce$, $\chainId$.
\fld{Steps.}
\begin{enumerate}[nosep,leftmargin=1.7em,topsep=2pt]
\item $\assert \initFlag \geq 1$, \; $\tokType \neq 0$, \; $\tokTypeB \neq 0$, \;
      $0 < \amtOut \leq 2^{64}-1$, \; $0 < \amtMax \leq 2^{64}-1$.
\item $\assert \coinIn.\mathsf{color} =
      \lnk{fn:tokenTypeOf}{\tokenTypeOf(\domSep(\tokType), \vaultAddr)}$ \; and \;
      $\coinIn.\mathsf{value} = \amtMax$.
\item $\evmTxOf \gets$ ``exact-output exchange of at most $\amtMax$ of $\tokType$ for exactly
      $\amtOut$ of $\tokTypeB$ at $\feeTier$, recipient  is $\vaultEvmAddr$, no price bound'',
      to $\routerAddr$, contract-fixed gas envelope.
\item $\keyPath \gets \keyPathLit$, as in $\cktStartWithdraw$.
\item $\req \gets (\vaultAddr\bc \sigNonce\bc \keyVer\bc \keyPath\bc \evmTxOf\bc \chainId)$; \;
      $\rid \gets \lnk{fn:reqIdOf}{\reqIdOf(\req)}$; \; $\assert \rid \notin \swapMap$.
\end{enumerate}
\MidNote{\textbf{Midnight operations: burn}, exactly as in \cktref{ckt:startwithdraw}:
$\mnhl{\code{receiveShielded}}(\coinIn)$ then
$\mnhl{\code{sendImmediateShielded}}(\coinIn, \burnAddr, \amtMax)$.}
\begin{OutList}
\oitem{\Dvault} $\sigNonce \mathrel{+}= 1$; \; $\swapMap[\rid] \gets \req$; \;
  \ADDEDINL{$\swapView[\rid] \gets \bigl(\lnk{fn:refundComm}{\refundComm(\usersk, \rid)},\;
  \tokType,\; \tokTypeB,\; \amtOut,\; \amtMax\bigr)$}.
\oitem{\Dshd} as \cktref{ckt:startwithdraw}, with $\amtMax$ in place of $\amt$: \;
  \bluedot $\coinCommit(\coinIn, \vaultAddr)$, \; \bluedot $\coinNull(\coinIn, \vaultAddr)$, \;
  \bluedot $\coinCommit(\coinBurnOut, \burnAddr)$, \;
  \bluedot $\coinNull(\coinFund, \coinsk)$ and $\coinCommit(\coinChange, \srcAcct)$.
\oitem{\stmt} $\bigl(\cktStartSwap,\; \req,\; \rid,\; \swapView[\rid],\; \Dshd,\;
  \notif(\vaultAddr, \rid, \swapMap),\; \comIO\bigr)$, where $\req$ exposes
  $\tokType, \tokTypeB, \feeTier, \amtOut, \amtMax$.
  \ADDEDINL{The view publishes $\tokType$, $\tokTypeB$, $\amtOut$ and $\amtMax$.}
\oitem{\wit} $(\usersk,\; \burnNonce,\; \coinsk)$.
\oitem{\Rel} as \cktref{ckt:startwithdraw}, over $\swapView$ and $\amtMax$.
\end{OutList}
\end{CktBox}


\begin{CktBox}{ckt:completeswap}{$\cktCompleteSwap$}
\fld{Arguments.} $\rid$, $\attSig$, $\evmOut$ of width $8$, $\freshrn{\mintNonce}$,
\ADDEDINL{$\freshrn{\changeNonce}$}.
\fld{Witness.} $\usersk$, $\mintNonce$, \ADDEDINL{$\changeNonce$}.
\fld{Reads.} $\initFlag$, $\mpcRespKey$, \ADDEDINL{$\swapView[\rid]$}, $\swapMap$.
\fld{Steps.}
\begin{enumerate}[nosep,leftmargin=1.7em,topsep=2pt]
\item $\assert \initFlag \geq 1$; \; \ADDEDINL{$\assert \changeNonce \neq \mintNonce$}.
\item $\assert \EcdsaVf\bigl(\lnk{fn:attDigest}{\attDigest(\rid, \evmOut)}\bc \attSig\bc
      \mpcRespKey\bigr)$.
\item \ADDEDINL{$\assert \rid \in \swapView$}; \; remove $\swapMap[\rid]$.
\item \ADDEDINL{$\assert \lnk{fn:refundComm}{\refundComm(\usersk, \rid)} = \swapView[\rid].\fComm$}
      (authenticating the swapper); \; remove $\swapView[\rid]$.
\item \ADDEDINL{$\amtOut \gets \swapView[\rid].\fAmtOut$,}
      \ADDEDINL{$\tokTypeB \gets \swapView[\rid].\fTokOut$,}
      \ADDEDINL{$\amtMax \gets \swapView[\rid].\fAmtMax$,}
      \ADDEDINL{$\tokType \gets \swapView[\rid].\fTokIn$}; \;
      $\amtIn \gets$ the attested value in $\evmOut$. \;
      $\amtRem \gets \amtMax - \amtIn$.
\end{enumerate}
\MidNote{\textbf{Midnight operations: two mints.} $\mintShieldedToken$ is called twice, both times
to the caller's own shielded account:
$\coinMint = \bigl(\mintNonce,\; \tokenTypeOf(\domSep(\tokTypeB), \vaultAddr),\; \amtOut\bigr)$ and
$\coinRem = \bigl(\changeNonce,\; \tokenTypeOf(\domSep(\tokType), \vaultAddr),\; \amtRem\bigr)$.
The second runs even when $\amtRem = 0$.}
\begin{OutList}
\oitem{\Dvault} $\swapMap[\rid]$ removed; \ADDEDINL{$\swapView[\rid]$ removed}.
\oitem{\Dshd} \bluedot mint record $\bigl(\domSep(\tokTypeB),\, \amtOut\bigr)$; \;
  \bluedot mint record $\bigl(\domSep(\tokType),\, \amtRem\bigr)$; \;
  \bluedot $\coinCommit(\coinMint, \text{caller})$; \;
  \bluedot $\coinCommit(\coinRem, \text{caller})$.
\oitem{\stmt} $\bigl(\cktCompleteSwap,\; \rid,\; \attSig,\; \amtIn,\; \Dvault,\; \Dshd\bigr)$.
\oitem{\wit} $(\usersk,\; \mintNonce,\; \ADDEDINL{\changeNonce})$.
\oitem{\Rel} $\refundComm(\usersk, \rid)$ $= \swapView[\rid].\fComm$ \;$\wedge$\;
  $\EcdsaVf(\attDigest(\rid,\evmOut)$, $\attSig$, $\mpcRespKey)$ \;$\wedge$\;
  $\changeNonce$ $\neq \mintNonce$ \;$\wedge$\;
  the first commitment opens to $(\mintNonce, \cdot, \amtOut, \text{caller})$ \;$\wedge$\;
  the second opens to $(\changeNonce, \cdot, \amtRem, \text{caller})$.
\end{OutList}
\end{CktBox}

\begin{CktBox}{ckt:refundswap}{$\cktRefundSwap$ }
\fld{Arguments.} $\rid$, $\attSig$, $\evmOut$ of with 5, $\mintNonce$.
\fld{Witness.} $\usersk$.
\fld{Reads.} $\initFlag$, $\mpcRespKey$, $\swapView[\rid]$, $\swapMap$.
\fld{Steps.}
\begin{enumerate}[nosep,leftmargin=1.7em,topsep=2pt]
\item Run \hyperref[ckt:failgate]{$\cktFailGate(\rid, \attSig, \evmOut)$}.
\item $\assert \rid \in \swapView$.
\item $\assert \lnk{fn:refundComm}{\refundComm(\usersk, \rid)} = \swapView[\rid].\fComm$ -
      \textbf{the swapper gate}.
\item Remove $\swapMap[\rid]$ and $\swapView[\rid]$.
\end{enumerate}
\MidNote{\textbf{Midnight operation.} Mint, to the caller's own shielded account,
$\coinMint = \bigl(\mintNonce,\; \mathsf{color},\;
\swapView[\rid].\fAmtMax\bigr).$}
where $\mathsf{color}= \lnk{fn:tokenTypeOf}{\tokenTypeOf(\domSep(\swapView[\rid].\fTokIn),
\vaultAddr)}$.
\begin{OutList}
\oitem{\Dvault} $\swapMap[\rid]$ removed; $\swapView[\rid]$ removed.
\oitem{\Dshd} \bluedot mint record
  $\bigl(\domSep(\swapView[\rid].\fTokIn),\, \swapView[\rid].\fAmtMax\bigr)$; \;
  \bluedot $\coinCommit(\coinMint, \text{caller})$.
\oitem{\stmt} $\bigl(\cktRefundSwap,\; \rid,\; \attSig,\; \Dvault,\; \Dshd\bigr)$.
\oitem{\wit} $(\usersk,\; \mintNonce)$.
\oitem{\Rel} $\refundComm(\usersk, \rid)$ $= \swapView[\rid].\fComm$ \;$\wedge$\; the relation of
  $\cktFailGate$ \;$\wedge$\; the published commitment opens to
  $(\mintNonce, \cdot, \swapView[\rid].\fAmtMax, \text{caller})$.
\end{OutList}
\end{CktBox}


\begin{PubBox}{What $\Swap$ publishes}
\textbf{On $\Lpub$:} the exchange in the clear: both tokens/colors and amounts, $\feeTier$, the recipient $\vaultEvmAddr$, $\amtMax$, $\amtOut$, and the resulting movements through the pool.

\smallskip
\textbf{On $\Lshd$:} as $\Withdraw$, over $\swapMap$ and $\swapMarkerMap$, with two mint
records at settlement. Mints do not leak any information about the payer and receiver.

\smallskip
\textbf{The executed price is public.} $\amtMax \in \req$ and $\amtIn \in \stmt_{\cktCompleteSwap}$
outright, and $\amtChg$ is published as a mint record as well. With both tokens, $\feeTier$ and the block position also public,every vault swap is a
fully transparent trade.  No address -  besides the vault address  - is involved in the trade.
\end{PubBox}


\subsection{\texorpdfstring{$\Supply$ and $\Redeem$}{Supply and Redeem}}
\label{ssec:lendinst}

The lending flow $\Supply$ and $\Redeem$ moves value between an ERC20 and its interest-bearing wrapper, both pinned at
$\Setup$ (\S\ref{ssec:state}), the token being lent is
at address $\underAddr$, the receipt token is at address $\wrapAddr$.
$\Supply$ and $\Redeem$ are both a $\Swap$.
$\Supply$ burns a coin of the underlying
color and mints one of the wrapper color; $\Redeem$ does the reverse.  They differ from
$\Swap$ in  that  \textbf{the minted amount is read from the attestation, not from
the request}. In $\Swap$, $\amtOut$ is a field of $\req$ that the caller fixed before
anything was burned; in $\Supply$ and $\Redeem$ the corresponding quantity is whatever the
lending market returned, and it reaches the circuit only through $\evmOut$.


Throughout the paper we use:
\begin{align*}
  \colrU &\;=\; \lnk{fn:tokenTypeOf}{\tokenTypeOf}\bigl(\domSep(\underAddr),\, \vaultAddr\bigr), \\
  \colrW &\;=\; \lnk{fn:tokenTypeOf}{\tokenTypeOf}\bigl(\domSep(\wrapAddr),\, \vaultAddr\bigr)
\end{align*}
for the two shielded colors involved: $\colrU$ the color of the lent ERC20, $\colrW$ the color of
the wrapper.

\subsubsection{Supply}

\begin{ProcBox}{$\Pi.\Supply^{\Lpub, \Lshd}\bigl(\shdAcct\bc \tokType\bc \amt\bc \auxx \,;\,
      \privIn\bigr)$}
\label{box:supply}
\sloppy
\emph{Private input.} $\privIn$ holds $\usersk$; the secret key $\coinsk$ associated to the
\zkswap\ account $\srcAcct$; a fresh \freshrn{$\burnNonce$}, used at Step~1; and a fresh
\freshrn{$\mintNonce$}, used at Step~3 on success or at Step~4 on failure.\\
\emph{Public arguments.} $\tokType = \underAddr$, read from pinned state rather than supplied;
$\auxx = (\evmNonce, \keyVer)$ only. The gas envelope is contract-fixed, because the
\emph{vault}'s own EVM account pays for this transaction. Requires $\cktApproveWrap$
(\cktref{ckt:approvestata}) to have been run once, so that the wrapper may pull the underlying.

\begin{Steps}
\item \textbf{Request call: $\cktStartSupply$} (\cktref{ckt:startsupply}). The caller hands the contract a
      coin descriptor $\coinIn = (\burnNonce,\, \colrU,\, \amt)$; the circuit burns it, files
      ``$\text{deposit}(\amt, \vaultEvmAddr)$ on $\wrapAddr$'' under $\rid$ in $\supplyMap$ with
      $\keyPath = \keyPathLit$, and pins
      $\supplyMarkerMap[\rid] \gets \lnk{fn:refundComm}{\refundComm(\usersk, \rid)}$. The caller address does not appear in the transaction.
\item \textbf{MPC.} The committee signs with the key derived at $\keyPath = \keyPathLit$, so from
      $\vaultEvmAddr$, the pool's own account; the relayer broadcasts; the committee observes the
      execution, reads the returned share count against $\schemaIn$, re-serialises it against
      $\schemaOut$, and posts $\attSig$ over
      $\lnk{fn:attDigest}{\attDigest(\rid, \evmOut)}$.
\item \textbf{Settlement call on success: $\cktCompleteSupply$} (\cktref{ckt:completesupply}).
      Supplier-gated; mints \emph{one} coin, of color $\colrW$ and value $\amtShares$ read from
      $\evmOut$.
\item \textbf{Settlement call on failure:  execute $\cktRefundSupply$}   (\cktref{ckt:refundsupply}):  Re-mints the \emph{full} $\amt$ of color $\colrU$ - the deposit never executed, so nothing was spent.
\end{Steps}

\medskip
Outcome: $\succeeded$ via Step~3, $\failed$ via Step~4. Private output: $\usersk$, and the new coins minted.
\end{ProcBox}


\begin{CktBox}{ckt:approvestata}{$\cktApproveStata$}
\fld{Arguments.} $\evmNonce$, $\keyVer$.
\fld{Witness.} none.
\fld{Reads.} $\initFlag$, $\underAddr$, $\wrapAddr$, $\vaultAddr$, $\sigNonce$, $\chainId$.
\fld{Steps.}
\begin{enumerate}[nosep,leftmargin=1.7em,topsep=2pt]
\item $\assert \initFlag \geq 1$.
\item $\evmTxOf \gets$ ``approve $\wrapAddr$ for an unlimited amount'', sent to $\underAddr$,
      contract-fixed gas envelope.
\item $\keyPath \gets \keyPathLit$; \;
      $\req \gets (\vaultAddr\bc \sigNonce\bc \keyVer\bc \keyPath\bc \evmTxOf\bc \chainId)$; \;
      $\rid \gets \reqIdOf(\req)$; \; $\assert \rid \notin \reqMap$.
\end{enumerate}
\begin{OutList}
\oitem{\Dvault} $\sigNonce \mathrel{+}= 1$; \; $\reqMap[\rid] \gets \req$.
\oitem{\Dshd} $\emptyset$.
\oitem{\stmt} $\bigl(\cktApproveStata,\; \req,\; \rid,\; \notif(\vaultAddr, \rid, \reqMap),\;
  \comIO\bigr)$.
\oitem{\wit} $\bot$.
\oitem{\Rel} the request is well formed.
\end{OutList}
\end{CktBox}


\begin{CktBox}{ckt:startsupply}{$\srcStartSupply: \cktStartSupply$}
\fld{Arguments.} $\evmNonce$, $\keyVer$, $\amt$, and
$\coinIn = (\burnNonce, \mathsf{color}, \mathsf{value})$.
\fld{Witness.} $\usersk$ (and $\coinsk$ at the wallet layer).
\fld{Reads.} $\initFlag$, $\underAddr$, $\wrapAddr$, $\vaultEvmAddr$, $\vaultAddr$, $\sigNonce$,
$\chainId$.
\fld{Steps.}
\begin{enumerate}[nosep,leftmargin=1.7em,topsep=2pt]
\item $\assert \initFlag \geq 1$, \; $0 < \amt \leq 2^{64}-1$.
\item $\assert \coinIn.\mathsf{color} = \colrU$ \; and \; $\coinIn.\mathsf{value} = \amt$.
\item $\evmTxOf \gets$ ``$\srcDepositSel$ for $\amt$, recipient $\vaultEvmAddr$'', sent to
      $\wrapAddr$, contract-fixed gas envelope.
\item $\keyPath \gets \keyPathLit$.
\item $\req \gets (\vaultAddr\bc \sigNonce\bc \keyVer\bc \keyPath\bc \evmTxOf\bc \chainId)$; \;
      $\rid \gets \lnk{fn:reqIdOf}{\reqIdOf(\req)}$; \; $\assert \rid \notin \supplyMap$.
\end{enumerate}
\MidNote{\textbf{Midnight operations: burn.}
$\mnhl{\code{receiveShielded}}(\coinIn)$ then
$\mnhl{\code{sendImmediateShielded}}(\coinIn, \burnAddr, \amt)$.}
\begin{OutList}
\oitem{\Dvault} $\sigNonce \mathrel{+}= 1$; \; $\supplyMap[\rid] \gets \req$; \;
  \ADDEDINL{$\supplyView[\rid] \gets \bigl(\lnk{fn:refundComm}{\refundComm(\usersk, \rid)},\;
  \amt\bigr)$}.
\oitem{\Dshd} \bluedot $\coinCommit(\coinIn, \vaultAddr)$, \;
  \bluedot $\coinNull(\coinIn, \vaultAddr)$, \;
  \bluedot $\coinCommit(\coinBurnOut, \burnAddr)$, \;
  \bluedot $\{\coinNull(\coinFund, \coinsk)\}$ and $\coinCommit(\coinChange, \srcAcct)$.
\oitem{\stmt} $\bigl(\cktStartSupply,\; \req,\; \rid,\; \supplyView[\rid],\; \Dshd,\;
  \notif(\vaultAddr, \rid, \supplyMap),\; \comIO\bigr)$.
\oitem{\wit} $(\usersk,\; \burnNonce,\; \coinsk)$.
\oitem{\Rel} $\coinIn.\mathsf{color} = \colrU$ \;$\wedge$\; $\coinIn.\mathsf{value} = \amt$
  \;$\wedge$\; $\supplyView[\rid].\fComm = \refundComm(\usersk, \rid)$.
\end{OutList}
\end{CktBox}


\begin{CktBox}{ckt:completesupply}{$\cktCompleteSupply$}
\fld{Arguments.} $\rid$, $\attSig$, $\evmOut$ of width $8$, $\freshrn{\mintNonce}$.
\fld{Witness.} $\usersk$, $\mintNonce$.
\fld{Reads.} $\initFlag$, $\mpcRespKey$, \ADDEDINL{$\supplyView[\rid]$}, $\supplyMap$, $\wrapAddr$.
\fld{Steps.}
\begin{enumerate}[nosep,leftmargin=1.7em,topsep=2pt]
\item $\assert \initFlag \geq 1$ and
      $\assert \EcdsaVf\bigl(\lnk{fn:attDigest}{\attDigest(\rid, \evmOut)}\bc \attSig\bc
      \mpcRespKey\bigr)$.
\item $\assert \rid \in \supplyMap$; \; remove $\supplyMap[\rid]$.
\item \ADDEDINL{$\assert \lnk{fn:refundComm}{\refundComm(\usersk, \rid)} =
      \supplyView[\rid].\fComm$} (authentication \textbf{supplier-only}); \; remove $\supplyView[\rid]$.
\item $\amtShares \gets$ the attested value in $\evmOut$, released publicly.
\end{enumerate}
\MidNote{\textbf{Midnight operation.} Mint
$\coinMintW = \bigl(\freshrn{\mintNonce},\; \colrW,\; \amtShares\bigr)$
to $\usrTag{\coinpk \text{ of } \setlAcct}$.}
\begin{OutList}
\oitem{\Dvault} $\supplyMap[\rid]$ removed; \ADDEDINL{$\supplyView[\rid]$ removed}.
\oitem{\Dshd} \bluedot mint record $\bigl(\domSep(\wrapAddr),\, \amtShares\bigr)$; \;
  \bluedot $\coinCommit(\coinMintW$, $\coinpk$ of $\shdAcct)$.
\oitem{\stmt} $\bigl(\cktCompleteSupply,\; \rid,\; \attSig,\; \amtShares,\; \Dvault,\;
  \Dshd\bigr)$.
\oitem{\wit} $(\usersk,\; \mintNonce)$.
\oitem{\Rel} $\refundComm(\usersk, \rid)$ $= \supplyView[\rid].\fComm$ \;$\wedge$\;
  $\EcdsaVf(\attDigest(\rid,\evmOut)$, $\attSig$, $\mpcRespKey)$ \;$\wedge$\;
  the commitment opens to
  $\bigl(\mintNonce,\, \colrW,\, \amtShares,\, \coinpk\bigr)$.
\end{OutList}
\end{CktBox}


\begin{CktBox}{ckt:refundsupply}{$\cktRefundSupply$}
\fld{Arguments.} $\rid$, $\attSig$, $\evmOut$ of with 5, $\mintNonce$.
\fld{Witness.} $\usersk$.
\fld{Reads.} $\initFlag$, $\mpcRespKey$, $\supplyView[\rid]$, $\supplyMap$.
\fld{Steps.}
\begin{enumerate}[nosep,leftmargin=1.7em,topsep=2pt]
\item Run \hyperref[ckt:failgate]{$\cktFailGate(\rid, \attSig, \evmOut)$}.
\item $\assert \rid \in \supplyView$.
\item $\assert \lnk{fn:refundComm}{\refundComm(\usersk, \rid)} = \supplyView[\rid].\fComm$ - (authentication:
      \textbf{supplier-only}).
\item Remove $\supplyMap[\rid]$ and $\supplyView[\rid]$.
\end{enumerate}
\MidNote{\textbf{Midnight operation.} Mint
$\coinMint = \bigl(\freshrn{\mintNonce},\; \colrU,\; \supplyView[\rid].\fAmt\bigr)$
to $\usrTag{\coinpk \text{ of } \setlAcct}$.}
\begin{OutList}
\oitem{\Dvault} $\supplyMap[\rid]$ removed; $\supplyView[\rid]$ removed.
\oitem{\Dshd} \bluedot mint record
  $\bigl(\domSep(\underAddr),\, \supplyView[\rid].\fAmt\bigr)$; \;
  \bluedot $\coinCommit\bigl(\coinMint,\, \usrTag{\coinpk \text{ of } \setlAcct}\bigr)$.
\oitem{\stmt} $\bigl(\cktRefundSupply,\; \rid,\; \attSig,\; \Dvault,\; \Dshd\bigr)$.
\oitem{\wit} $(\usersk,\; \mintNonce)$.
\oitem{\Rel} $\refundComm(\usersk, \rid)$ $= \supplyView[\rid].\fComm$ \;$\wedge$\; the relation of
  $\cktFailGate$ \;$\wedge$\; the commitment opens to
  $\bigl(\mintNonce,\, \colrU,\, \supplyView[\rid].\fAmt,\, \usrTag{\coinpk \text{ of }
  \shdAcct}\bigr)$.
\end{OutList}
\end{CktBox}

\begin{PubBox}{What $\Supply$ publishes}
\textbf{On $\Lpub$:} the supply in the clear: the wrapper contract called, the amount of the
underlying supplied, the receiver $\vaultEvmAddr$, and the share count returned. The vault's own account is the sender, so \emph{nothing tied to $\srcAcct$ appears}.

\smallskip
\textbf{On $\Lshd$:} $\stmt_{\cktSupply}$, the MPC signatures, and the settlement statement.
Of these, the components that involve the caller at all are: the marker
$\refundComm(\usersk, \rid)$; the coins burned to the contract: the nullifier and the new coin commitment destined to $\burnAddr$ (the private and random nonce $\burnNonce$ protects the identity of the caller of $\cktStartSupply$), the new private coins minted in the $\cktCompleteSupply/\cktRefundSupply$ destined  to $\shdAcct$ (the private and random nonce $\mintNonce$ protects the identity of the recipient).

\smallskip
\textbf{The exchange rate is public.} $\amt$ is a field of $\req$ and $\amtShares$ is
released by the settlement call and published again as a mint record, so
$\amt / \amtShares$. This is
information about $\Lpub$, not about the caller, and it is identical for every supplier in the same
block.

\smallskip
\textbf{Hidden:} which shielded account funded the call; which shielded account received the shares.\end{PubBox}


\subsubsection{Redeem}

\begin{ProcBox}{$\Pi.\Redeem^{\Lpub, \Lshd}\bigl(\shdAcct\bc \wrapOf{\tokType}\bc \amtShares\bc
      \auxx \,;\, \privIn\bigr)$}
\label{box:redeem}
\sloppy
\emph{Private input.} As $\Supply$: $\usersk$; $\coinsk$ of $\srcAcct$; a fresh
\freshrn{$\burnNonce$} at Step~1; and a fresh \freshrn{$\mintNonce$} at Step~3 on success or Step~4
on failure.\\
\emph{Public arguments.} $\wrapOf{\tokType} = \wrapAddr$, read from pinned state;
$\auxx = (\evmNonce, \keyVer)$ only. No approval leg is needed: the vault redeems its own shares,
so no allowance to a third party is involved.

\begin{Steps}
\item \textbf{Request call: $\cktRedeem$} (\cktref{ckt:startredeem}). The caller hands the contract a coin
      descriptor $\coinIn = (\burnNonce,\, \colrW,\, \amtShares)$; the circuit burns it, files
      ``$\text{redeem}(\amtShares, \vaultEvmAddr, \vaultEvmAddr)$ on $\wrapAddr$'' under $\rid$ in
      $\redeemMap$ with $\keyPath = \keyPathLit$, and pins
      $\redeemMarkerMap[\rid] \gets \lnk{fn:refundComm}{\refundComm(\usersk, \rid)}$.
\item \textbf{MPC.} As in $\Supply$, signing from $\vaultEvmAddr$. The committee reports
      $\amtAssets$, the amount of the underlying the market returned --- principal plus accrued
      interest.
\item \textbf{Settlement call on success: $\cktCompleteRedeem$} (\cktref{ckt:completeredeem}).
      Redeemer-gated; mints \emph{one} coin, of colour $\colrU$ and value $\amtAssets$.
\item \textbf{Settlement call on failure: $\cktRefundRedeem$}
      (\cktref{ckt:refundredeem}). Re-mints the \emph{full} $\amtShares$ of color $\colrW$.
\end{Steps}

\medskip
Outcome: $\succeeded$ via Step~3, $\failed$ via Step~4, $\pending$ otherwise. Private output:
$\usersk$.
\end{ProcBox}

% \begin{CktBox}{ckt:redeem}{$\cktRedeem$ --- the request call of $\Redeem$}
% \fld{Arguments.} $\evmNonce$, $\keyVer$, $\amtShares$, and the coin to surrender
% $\coinIn = (\burnNonce,\; \colr,\; \mathsf{value})$.
% \fld{Witness.} $\usersk$; and, at the wallet layer, $\coinsk$ of $\srcAcct$ and the openings of the
% coins the wallet spends.
% \fld{Reads.} $\initFlag$, $\wrapAddr$, $\vaultEvmAddr$, $\vaultAddr$, $\sigNonce$, $\chainId$.
% \fld{Steps.}
% \begin{enumerate}[nosep,leftmargin=1.7em,topsep=2pt]
% \item $\assert \initFlag \geq 1$, \; $0 < \amtShares \leq 2^{64}-1$. \textbf{This bounds the
%       \emph{input}, and only the input.} The refund path re-mints $\amtShares$, so the bound makes
%       the refund always representable; it says nothing about $\amtAssets$, which is what the success
%       path mints. That asymmetry is the subject of Warning~\ref{warn:redeembound}.
% \item $\assert \coinIn.\colr = \colrW$ \; and \; $\coinIn.\mathsf{value} = \amtShares$. As in
%       $\cktSupply$, both compared values are public and $\burnNonce$ is unconstrained.
% \item $\evmTxOf \gets$ ``$\srcRedeemSel$: redeem $\amtShares$ shares, receiver $\vaultEvmAddr$,
%       owner $\vaultEvmAddr$'', sent to $\wrapAddr$, on $\chainId$, no native value,
%       \textbf{contract-fixed} gas envelope. \emph{Both the receiver and the owner are read from
%       pinned state}: the shares being redeemed are the pool's, and the assets return to the pool.
% \item $\keyPath \gets \keyPathLit$.
% \item $\req \gets (\vaultAddr\bc \sigNonce\bc \keyVer\bc \keyPath\bc \evmTxOf\bc \chainId\bc
%       \schemaIn\bc \schemaOut)$, with the two contract-fixed redeem schemas; \;
%       $\rid \gets \lnk{fn:reqIdOf}{\reqIdOf(\req)}$; \; $\assert \rid \notin \redeemMap$.
% \end{enumerate}

% \MidNote{\textbf{Midnight operations: burn redeemed coin}, as in \cktref{ckt:supply}:
% $\mnhl{\opReceive}(\coinIn)$ then $\mnhl{\opSendImm}(\coinIn, \burnAddr, \amtShares)$.}

% \smallskip
% \noindent\textbf{The four coins this call touches.} Identical in structure to the table of
% \cktref{ckt:supply}, with $\colrW$ in place of $\colrU$ and $\amtShares$ in place of $\amt$:
% $\coinIn = (\burnNonce,\, \colrW,\, \amtShares)$;
% $\coinFund_i$ of colour $\colrW$ summing to $V \geq \amtShares$;
% $\coinChange = (\cdot,\, \colrW,\, V - \amtShares)$;
% $\coinBurnOut = \bigl(\evolveNonce(\burnNonce),\, \colrW,\, \amtShares\bigr)$.

% \begin{OutList}
% \oitem{\Dvault} $\sigNonce \mathrel{+}= 1$; \; $\redeemMap[\rid] \gets \req$; \;
%   $\redeemMarkerMap[\rid] \gets \lnk{fn:refundComm}{\refundComm(\usersk, \rid)}$.
% \oitem{\Dshd} exactly the five entries of \cktref{ckt:supply}, with $\colrW$ and $\amtShares$
%   substituted: \; \bluedot $\{\coinNull(\coinFund_i, \usrTag{\coinsk})\}_i$, \;
%   \bluedot $\coinNull(\coinIn, \ctrTag{\vaultAddr})$, \;
%   \bluedot $\coinCommit(\coinIn, \ctrTag{\vaultAddr})$, \;
%   \bluedot $\coinCommit(\coinBurnOut, \burnAddr)$, \;
%   \bluedot $\coinCommit(\coinChange, \usrTag{\coinpk})$.
% \oitem{\stmt} $\bigl(\cktRedeem,$ $\req,$ $\rid,$ $\refundComm(\usersk, \rid),$ $\Dshd,$
%   $\notif(\vaultAddr, \rid, \redeemMap),$ $\comIO\bigr)$, where $\req$ exposes
%   $\wrapAddr, \amtShares, \vaultEvmAddr, \evmNonce, \sigNonce$ and $\keyPath = \keyPathLit$. The
%   notification \textbf{names $\redeemMap$}.
% \oitem{\wit} $\bigl(\usersk,\; \burnNonce,\; \coinsk,\; \{\coinFund_i\}_i\bigr)$.
% \oitem{\Rel} $\redeemMarkerMap[\rid] = \refundComm(\usersk, \rid)$ \;$\wedge$\;
%   $\coinIn = (\burnNonce, \colrW, \amtShares)$ opens both
%   $\coinCommit(\coinIn, \ctrTag{\vaultAddr})$ and $\coinNull(\coinIn, \ctrTag{\vaultAddr})$
%   \;$\wedge$\; each $\coinFund_i$ opens its published nullifier under $\coinsk$ and is unspent
%   \;$\wedge$\; the colour balances.
% \end{OutList}
% \end{CktBox}

% \ADDEDMARK{Current source: \srcStartRedeem. Changes from Circuit~C15: the pending marker is now
% $\redeemView[\rid]$, a record holding the caller commitment and the share count, in place of
% $\redeemMarkerMap[\rid]$, which held the commitment alone.}
\begin{CktBox}{ckt:startredeem}{$\cktStartRedeem$ : (\srcStartRedeem)}
\fld{Arguments.} $\evmNonce$, $\keyVer$, $\amtShares$, and
$\coinIn = (\burnNonce, \mathsf{color}, \mathsf{value})$.
\fld{Witness.} $\usersk$ (and $\coinsk$ at the wallet layer).
\fld{Reads.} $\initFlag$, $\wrapAddr$, $\vaultEvmAddr$, $\vaultAddr$, $\sigNonce$, $\chainId$.
\fld{Steps.}
\begin{enumerate}[nosep,leftmargin=1.7em,topsep=2pt]
\item $\assert \initFlag \geq 1$, \; $0 < \amtShares \leq 2^{64}-1$.
\item $\assert \coinIn.\mathsf{color} = \colrW$ \; and \; $\coinIn.\mathsf{value} = \amtShares$.
\item $\evmTxOf \gets$ ``$\srcRedeemSel$ for $\amtShares$, receiver $\vaultEvmAddr$, owner
      $\vaultEvmAddr$'', sent to $\wrapAddr$, contract-fixed gas envelope.
\item $\keyPath \gets \keyPathLit$.
\item $\req \gets (\vaultAddr\bc \sigNonce\bc \keyVer\bc \keyPath\bc \evmTxOf\bc \chainId)$; \;
      $\rid \gets \lnk{fn:reqIdOf}{\reqIdOf(\req)}$; \; $\assert \rid \notin \redeemMap$.
\end{enumerate}
\MidNote{\textbf{Midnight operations: burn.}
$\mnhl{\code{receiveShielded}}(\coinIn)$ then
$\mnhl{\code{sendImmediateShielded}}(\coinIn, \burnAddr, \amtShares)$.}
\begin{OutList}
\oitem{\Dvault} $\sigNonce \mathrel{+}= 1$; \; $\redeemMap[\rid] \gets \req$; \;
  \ADDEDINL{$\redeemView[\rid] \gets \bigl(\lnk{fn:refundComm}{\refundComm(\usersk, \rid)},\;
  \amtShares\bigr)$}.
\oitem{\Dshd} \bluedot $\coinCommit(\coinIn, \vaultAddr)$, \;
  \bluedot $\coinNull(\coinIn, \vaultAddr)$, \;
  \bluedot $\coinCommit(\coinBurnOut, \burnAddr)$, \;
  \bluedot $\{\coinNull(\coinFund_i, \coinsk)\}_i$ and $\coinCommit(\coinChange, \srcAcct)$.
\oitem{\stmt} $\bigl(\cktStartRedeem,\; \req,\; \rid,\; \redeemView[\rid],\; \Dshd,\;
  \notif(\vaultAddr, \rid, \redeemMap),\; \comIO\bigr)$.
\oitem{\wit} $(\usersk,\; \burnNonce,\; \coinsk)$.
\oitem{\Rel} $\coinIn.\mathsf{color} = \colrW$ \;$\wedge$\; $\coinIn.\mathsf{value} = \amtShares$
  \;$\wedge$\; $\redeemView[\rid].\fComm = \refundComm(\usersk, \rid)$.
\end{OutList}
\end{CktBox}


\begin{CktBox}{ckt:completeredeem}{$\cktCompleteRedeem$}
\fld{Arguments.} $\rid$, $\attSig$, $\evmOut$ of width $8$, $\freshrn{\mintNonce}$.
\fld{Witness.} $\usersk$, $\mintNonce$.
\fld{Reads.} $\initFlag$, $\mpcRespKey$, \ADDEDINL{$\redeemView[\rid]$}, $\redeemMap$,
$\underAddr$.
\fld{Steps.}
\begin{enumerate}[nosep,leftmargin=1.7em,topsep=2pt]
\item $\assert \initFlag \geq 1$ and
      $\assert \EcdsaVf\bigl(\lnk{fn:attDigest}{\attDigest(\rid, \evmOut)}\bc \attSig\bc
      \mpcRespKey\bigr)$.
\item $\assert \rid \in \redeemMap$; \; remove $\redeemMap[\rid]$.
\item \ADDEDINL{$\assert \lnk{fn:refundComm}{\refundComm(\usersk, \rid)} =
      \redeemView[\rid].\fComm$} --- \textbf{redeemer-only}; \; remove $\redeemView[\rid]$.
\item $\amtAssets \gets$ the attested value in $\evmOut$, \textbf{released publicly}.
\end{enumerate}
\MidNote{\textbf{Midnight operation.} Mint
$\coinMintU = \bigl(\freshrn{\mintNonce},\; \colrU,\; \amtAssets\bigr)$
to $\coinpk$.}
\begin{OutList}
\oitem{\Dvault} $\redeemMap[\rid]$ removed; \ADDEDINL{$\redeemView[\rid]$ removed}.
\oitem{\Dshd} \bluedot mint record $\bigl(\domSep(\underAddr),\, \amtAssets\bigr)$; \;
  \bluedot $\coinCommit\bigl(\coinMintU,\, \coinpk\bigr)$.
\oitem{\stmt} $\bigl(\cktCompleteRedeem,\; \rid,\; \attSig,\; \amtAssets,\; \Dvault,\;
  \Dshd\bigr)$.
\oitem{\wit} $(\usersk,\; \mintNonce)$.
\oitem{\Rel} $\refundComm(\usersk, \rid)$ $= \redeemView[\rid].\fComm$ \;$\wedge$\;
  $\EcdsaVf(\attDigest(\rid,\evmOut)$, $\attSig$, $\mpcRespKey)$ \;$\wedge$\;
  the commitment opens to
  $\bigl(\mintNonce,\, \colrU,\, \amtAssets,\, \coinpk\bigr)$.
\end{OutList}
\end{CktBox}


\begin{CktBox}{ckt:refundredeem}{$\cktRefundRedeem$}
\fld{Arguments.} $\rid$, $\attSig$, $\evmOut$ of width $5$, $\mintNonce$.
\fld{Witness.} $\usersk$.
\fld{Reads.} $\initFlag$, $\mpcRespKey$, $\redeemView[\rid]$, $\redeemMap$.
\fld{Steps.}
\begin{enumerate}[nosep,leftmargin=1.7em,topsep=2pt]
\item Run \hyperref[ckt:failgate]{$\cktFailGate(\rid, \attSig, \evmOut)$}.

\item $\assert \rid \in \redeemView$.
\item $\assert \lnk{fn:refundComm}{\refundComm(\usersk, \rid)} = \redeemView[\rid].\fComm$ (authenticating
      \textbf{redeemer-only}).
\item Remove $\redeemMap[\rid]$ and $\redeemView[\rid]$.
\end{enumerate}
\MidNote{\textbf{Midnight operation.} Mint
$\coinMint = \bigl(\freshrn{\mintNonce},\; \colrW,\; \redeemView[\rid].\fShares\bigr)$
to $\usrTag{\coinpk \text{ of } \setlAcct}$.}
\begin{OutList}
\oitem{\Dvault} $\redeemMap[\rid]$ removed; $\redeemView[\rid]$ removed.
\oitem{\Dshd} \bluedot mint record
  $\bigl(\domSep(\wrapAddr),\, \redeemView[\rid].\fShares\bigr)$; \;
  \bluedot $\coinCommit\bigl(\coinMint,\, \usrTag{\coinpk \text{ of } \setlAcct}\bigr)$.
\oitem{\stmt} $\bigl(\cktRefundRedeem,\; \rid,\; \attSig,\; \Dvault,\; \Dshd\bigr)$.
\oitem{\wit} $(\usersk,\; \mintNonce)$.
\oitem{\Rel} $\refundComm(\usersk, \rid)$ $= \redeemView[\rid].\fComm$ \;$\wedge$\; the relation of
  $\cktFailGate$ \;$\wedge$\; the commitment opens to
  $\bigl(\mintNonce,\, \colrW,\, \redeemView[\rid].\fShares,\, \usrTag{\coinpk \text{ of }
  \setlAcct}\bigr)$.
\end{OutList}
\end{CktBox}
\begin{PubBox}{What $\Redeem$ publishes}
\textbf{On $\Lpub$:} the redemption in the clear: the wrapper called, the share count burned,  and the amount of the underlying returned. The vault's account $\vaultEvmAddr$ is involved in all transactions on $\Lpub$ so nothing tied to $\srcAcct$ appears.

\smallskip
\textbf{On $\Lshd$:} as $\Supply$, over $\redeemMap$ and $\redeemMarkerMap$, with the two colors exchanged.

\smallskip
\textbf{The accrued yield is public.} $\amtShares$ is a field of $\req$ and $\amtAssets$ is released
by the settlement call, so the realised rate is recoverable, and with it the interest the position
earned.

\smallskip
\textbf{Hidden:} which shielded account funded the call; which received the assets.
\end{PubBox}
